% -----------------------------------------------------------------------------
\subsubsection{Liste des pièces (F1)}

La liste des pièces est l'écran d'accueil de l'application. Elle répond d'un coup d'œil à
trois questions~: quelle température fait-il dans chaque pièce, où chauffe-t-on en ce
moment, et avec quelle force. Chaque pièce occupe une ligne.

\begin{exigence}{TPB-PIE-001}{Une ligne par pièce}[
    justification = {Connaître la température de chaque pièce sans rien ouvrir.},
    priorite      = {Essentielle},
    verification  = {D},
    phase         = {Pièces},
  ]
  L'application doit présenter, pour chaque pièce chauffée de la maison Tado de
  l'utilisateur, une ligne portant le nom de la pièce et sa température mesurée.
\end{exigence}

\begin{exigence}{TPB-PIE-002}{Chauffe en cours}[
    justification = {Savoir où le chauffage tourne réellement, et pas seulement où il est
                     demandé.},
    priorite      = {Essentielle},
    verification  = {T},
    phase         = {Pièces},
    note          = {Une pièce chauffe lorsque sa puissance de chauffe est supérieure à
                     zéro.},
  ]
  L'application doit indiquer, sur la ligne de chaque pièce, si elle est en train de
  chauffer.
\end{exigence}

\begin{exigence}{TPB-PIE-003}{Puissance de chauffe}[
    justification = {Une pièce qui chauffe à 15~\% et une autre à 100~\% ne se trouvent
                     pas dans la même situation.},
    priorite      = {Essentielle},
    verification  = {T},
    phase         = {Pièces},
  ]
  Lorsqu'une pièce chauffe, l'application doit afficher sur sa ligne la puissance de
  chauffe, en pourcentage de la puissance maximale.
\end{exigence}

\begin{exigence}{TPB-PIE-004}{Température demandée}[
    justification = {Comparer ce qu'il fait à ce que l'on veut obtenir.},
    priorite      = {Importante},
    verification  = {D},
    phase         = {Pièces},
  ]
  L'application doit afficher sur la ligne de chaque pièce la température que le
  chauffage cherche à atteindre.
\end{exigence}

\begin{exigence}{TPB-PIE-005}{Pièce forcée signalée}[
    justification = {Un forçage oublié coûte de l'énergie~; il doit se voir sans ouvrir la
                     pièce.},
    priorite      = {Importante},
    verification  = {D},
    phase         = {Pièces},
  ]
  Lorsque la température d'une pièce a été imposée en dehors de la programmation,
  l'application doit le signaler sur sa ligne.
\end{exigence}

\begin{exigence}{TPB-PIE-006}{Pièce sans chauffage actif}[
    justification = {Une pièce dont le chauffage est éteint n'est pas une pièce qui se
                     refroidit par erreur.},
    priorite      = {Importante},
    verification  = {D},
    phase         = {Pièces},
  ]
  Lorsque le chauffage d'une pièce est éteint, l'application doit l'indiquer sur sa
  ligne, à la place de la température demandée.
\end{exigence}

\begin{exigence}{TPB-PIE-007}{Mesure indisponible}[
    justification = {Une température inventée ou périmée est pire qu'une absence de
                     température.},
    priorite      = {Essentielle},
    verification  = {T},
    phase         = {Pièces},
  ]
  Lorsque la température d'une pièce n'est pas connue, l'application doit l'indiquer sur
  sa ligne, sans afficher de valeur.
\end{exigence}

\begin{exigence}{TPB-PIE-008}{Ordre des pièces}[
    justification = {Retrouver chaque pièce à la même place d'une consultation à l'autre.},
    priorite      = {Importante},
    verification  = {D},
    phase         = {Pièces},
    note          = {L'ordre retenu est celui de l'application mobile Tado.},
  ]
  L'application doit présenter les pièces dans un ordre qui ne change pas d'une
  consultation à l'autre.
\end{exigence}

\begin{exigence}{TPB-PIE-009}{Beaucoup de pièces}[
    justification = {Une maison compte plus de pièces qu'un écran de montre n'en montre.},
    priorite      = {Essentielle},
    verification  = {T},
    phase         = {Pièces},
    note          = {Maison d'essai~: six pièces (annexe, §\ref{sec:maison-essai}). Le
                     nombre minimal de pièces à gérer reste à fixer.},
  ]
  L'application doit permettre d'atteindre toutes les pièces de la maison, même lorsqu'elles
  ne tiennent pas toutes à l'écran.
\end{exigence}

\begin{exigence}{TPB-PIE-010}{Présence dans le titre}[
    justification = {Savoir en permanence si le chauffage se croit à la maison ou
                     absent.},
    priorite      = {Importante},
    verification  = {D},
    phase         = {Pièces},
  ]
  L'application doit afficher, avec la liste des pièces, l'état de présence de la maison
  (\emph{Présent} ou \emph{Absent}).
\end{exigence}

\begin{exigence}{TPB-PIE-011}{Retour à la pièce choisie}[
    justification = {Après avoir forcé le salon, on veut retrouver le salon, pas le haut de
                     la liste.},
    priorite      = {Souhaitable},
    verification  = {D},
    phase         = {Pièces},
  ]
  Lorsque l'utilisateur revient à la liste des pièces depuis l'une d'elles, l'application
  doit sélectionner cette pièce.
\end{exigence}

% -----------------------------------------------------------------------------
\subsubsection{Détail d'une pièce (F2)}

Le détail d'une pièce rassemble ce que la liste ne peut pas montrer faute de place~: le
programme en cours, l'humidité, ce qui se passera ensuite. C'est aussi la porte d'entrée
vers le pilotage de la pièce.

\begin{exigence}{TPB-DET-001}{Détail d'une pièce}[
    justification = {Comprendre pourquoi une pièce chauffe, ou ne chauffe pas.},
    priorite      = {Essentielle},
    verification  = {D},
    phase         = {Pièces},
  ]
  L'application doit permettre d'ouvrir le détail d'une pièce depuis la liste, et y
  afficher sa température mesurée, la température demandée, l'état et la puissance de
  chauffe.
\end{exigence}

\begin{exigence}{TPB-DET-002}{Mode en cours}[
    justification = {Distinguer le programme d'un forçage, et savoir jusqu'à quand il
                     dure.},
    priorite      = {Essentielle},
    verification  = {T},
    phase         = {Pièces},
    note          = {Exemples~: \enquote{Programme, jusqu'à 22:00}~; \enquote{Forcé, encore
                     40 min}~; \enquote{Jusqu'à nouvel ordre}.},
  ]
  L'application doit afficher, dans le détail d'une pièce, si la température demandée vient
  de la programmation ou d'un forçage, et jusqu'à quel moment elle reste en vigueur.
\end{exigence}

\begin{exigence}{TPB-DET-003}{Prochain changement}[
    justification = {Savoir ce que le chauffage fera ensuite, avant d'y toucher.},
    priorite      = {Importante},
    verification  = {T},
    phase         = {Pièces},
  ]
  Lorsque la programmation prévoit un changement, l'application doit afficher dans le
  détail d'une pièce l'heure de ce changement et la température qui suivra.
\end{exigence}

\begin{exigence}{TPB-DET-004}{Humidité}[
    justification = {L'humidité relative accompagne la température dans la lecture d'une
                     pièce.},
    priorite      = {Importante},
    verification  = {T},
    phase         = {Pièces},
  ]
  Lorsque la pièce est équipée d'un capteur d'humidité, l'application doit afficher
  l'humidité relative dans le détail de la pièce.
\end{exigence}

\begin{exigence}{TPB-DET-005}{Passer d'une pièce à l'autre}[
    justification = {Faire le tour de la maison sans repasser par la liste.},
    priorite      = {Importante},
    verification  = {D},
    phase         = {Pièces},
  ]
  L'application doit permettre de passer, depuis le détail d'une pièce, au détail de la
  pièce précédente ou suivante de la liste.
\end{exigence}

\begin{exigence}{TPB-DET-006}{Accès au pilotage}[
    justification = {On consulte une pièce pour la régler.},
    priorite      = {Essentielle},
    verification  = {D},
    phase         = {Pièces},
  ]
  L'application doit permettre d'ouvrir le pilotage d'une pièce depuis son détail.
\end{exigence}

\begin{exigence}{TPB-DET-007}{Fenêtre ouverte dans le détail}[
    justification = {Expliquer un chauffage coupé alors que la programmation en demande.},
    priorite      = {Importante},
    verification  = {T},
    phase         = {Pièces},
  ]
  Lorsqu'une fenêtre ouverte est détectée dans une pièce, l'application doit l'indiquer
  dans le détail de la pièce, avec depuis combien de temps le chauffage est coupé.
\end{exigence}
